\section{Generation--Verification Gap}

\subsection{Oracle coverage 是能力上界，不是可交付性能}

Oracle verifier 知道 ground truth，因此总能从候选中挑出正确答案。真实系统通常只能使用多数投票、reward model 或 LLM judge，这些方法与 oracle 之间可能存在巨大差距。

\videofigure{figures/verification-gap.jpg}{常见选择方法远低于候选集合的 oracle coverage。}{00:15:32--00:18:44}

定义：

$$
G(k)=C_{\mathrm{oracle}}(k)-A_{\mathrm{selector}}(k),
$$

\begin{itemize}
    \item $C_{\mathrm{oracle}}(k)$：候选集合中存在正确答案的比例；
    \item $A_{\mathrm{selector}}(k)$：真实选择器交付正确答案的比例；
    \item $G(k)$：generation--verification gap。
\end{itemize}

\subsection{为什么 majority vote 会失败}

多数投票假设正确答案是分布的主模态。但对困难问题，模型可能产生许多相似错误，而正确轨迹极少出现。此时增加采样反而让错误模态的频率估计更稳定。

\videofigure{figures/rare-correct.jpg}{正确候选可能位于输出分布的稀有尾部。}{00:18:44--00:19:44}

\begin{warningbox}{Self-consistency 依赖一个隐藏假设}
多数投票只有在“正确推理路径比任一错误模式更集中”时才可靠。若模型系统性偏向同一种错误，self-consistency 会自信地选择错误答案。候选间的一致性不是正确性的充分证据。
\end{warningbox}

\subsection{验证器应具备什么性质}

\begin{itemize}
    \item \textbf{Soundness}：不要把明显错误判为正确；
    \item \textbf{Recall}：不要漏掉候选中的稀有正确解；
    \item \textbf{Calibration}：分数差异应反映真实成功概率；
    \item \textbf{Robustness}：对风格、长度和表面自信不过度敏感；
    \item \textbf{Cost}：验证成本应显著低于重新生成或人工检查。
\end{itemize}

\begin{knowledgebox}{可验证领域的结构优势}
数学形式化证明和代码执行可把 verifier 外包给符号系统或真实环境。开放式写作、战略决策和科学假设则缺少便宜 ground truth，只能使用 learned reward model 或延迟反馈，因此更容易出现验证瓶颈。
\end{knowledgebox}

\subsection{本章小结}

Inference scaling 的端到端上限由 verifier 决定。生成器越强、采样越多，候选分布越复杂，选择器反而可能面临更难的识别问题；因此 robust verification 是下一讲的核心主题。

